% Generated draft: requires mathematical and source review.
\documentclass[11pt]{article}
\usepackage[margin=25mm]{geometry}
\usepackage{fontspec}
\setmainfont{Latin Modern Roman}
\usepackage{amsmath,amssymb,mathtools,unicode-math}
\setmathfont{Latin Modern Math}
\usepackage{enumitem,needspace,xurl,hyperref}
\setlength{\parindent}{0pt}
\setlength{\parskip}{5pt}
\setlength{\emergencystretch}{3em}
\newcommand{\N}{\mathbb{N}}
\newcommand{\Z}{\mathbb{Z}}
\newcommand{\Q}{\mathbb{Q}}
\newcommand{\R}{\mathbb{R}}
\newcommand{\C}{\mathbb{C}}
\newcommand{\F}{\mathbb{F}}
\newcommand{\A}{\mathbb{A}}
\newcommand{\PP}{\mathbb{P}}
\newcommand{\E}{\mathbb{E}}
\newcommand{\eps}{\varepsilon}
\newcommand{\dd}{\,\mathrm{d}}
\begin{document}
\section*{UnsolvedMath ID: 30001068}
\textbf{Source problem code:} catalogue entry
\subsection{Definitions and mathematical statement}
Let $M$ be a finite simple rank-three matroid on $E=\{1,\ldots,n\}$, with real weights $a_i$ satisfying $\sum_{i\in E}a_i\geq0$. For $A\subseteq E$ put $w(A)=\sum_{i\in A}a_i$. A line is a rank-two flat, and a point is nonnegative when $a_i\geq0$. Form the incidence graph whose point vertices are $P_+=\{i:a_i\geq0\}$ and whose line vertices are rank-two flats $L$ with $w(L)\geq0$, joining $i$ to $L$ when $i\in L$. The first question asks whether all vertices in $P_+$ are connected through such nonnegative lines. Equivalently, are any two such points joined by an alternating point-line path in this incidence graph?\par For the higher-rank question, let $M$ have rank $r$ with the corresponding real weights and nonnegative total. A flat $F$ is characterized by $\operatorname{rk}(F\cup\{e\})>\operatorname{rk}(F)$ for every $e\notin F$. Let $\widehat0=\operatorname{cl}(\varnothing)$ and $\widehat1=E$ be the bottom and top flats. The usual proper-part interpretation of the weighted flat poset is \[\mathcal L_+(M,a)=\{F\text{ a flat}:\widehat0<F<\widehat1,\ w(F)\geq0\}.\] Its order complex has these flats as vertices and finite chains of them as simplices. Is this complex $(r-3)$-connected, and is it shellable? For $r\geq3$, the connectivity assertion means path connectedness and vanishing of the homotopy groups through degree $r-3$, interpreted as path connectedness alone at degree zero. The source does not specify low-rank conventions or whether simplicity carries over to the generalization. It also does not define proper to exclude the bottom flat, which is essential: in the loopless case including the empty flat would make the order complex a cone. The pure shellability convention orders facets so that each new facet meets the preceding union in a pure codimension-one subcomplex; a nonpure convention or purity assumption is not specified. All inequalities remain non-strict, with zero weights allowed; no genericity hypothesis is inserted.\par The supplied 2026-08-21 literature assessment is mismatched: it discusses parametric assignment, planar rotation matching, and polynomial intervals/local minima, rather than weighted matroids. Its linked geometric-optimization reference is retained for provenance but cannot be treated as establishing progress or status for this displayed question. The registry labels the matroid record open; that label is not independently verified here.
\subsection{Short English statement}
Ask whether nonnegative points are connected through nonnegative lines, and whether the proper nonnegative-flat order complex is highly connected and shellable. The treatment of the bottom flat and shellability conventions need clarification; the supplied literature summary concerns a different matching problem.
\subsection{Sources}
\begin{itemize}\item Local UnsolvedMath registry, permanent ID 30001068, source problem catalogue entry, \emph{Connectivity of Nonnegative Flats in Weighted Matroids}, proposed year 2008; no individual proposer supplied.\item Supplied citation: \emph{Discrete Geometry} (2009), Oberwolfach Reports, DOI 10.4171/owr/2008/44, pp.\ 2477--2556: \url{https://doi.org/10.4171/owr/2008/44}. Supplied literature label: OWR Report 44.\item Supplied literature label: Tessellations for geometric optimization, \url{https://refubium.fu-berlin.de/bitstream/fub188/419/1/final.pdf}. This is retained despite the supplied assessment's mismatch with the matroid question.\item Literature assessment date: 2026-08-21. The status and progress claims are supplied dated curation, not independently verified here; the assessment discusses planar rotation matching rather than the displayed weighted-matroid problem.\end{itemize}
\end{document}
